% from main.tex
%%%%%%%%%%  scan idx 116  |  folio 103  %%%%%%%%%%
implication \((\mathrm{ii})\Rightarrow(\mathrm{i})\) résulte aussitôt de (2.4) : la flèche \(0\to F\) est un 1-quasi-isomorphisme, on en déduit (quitte à localiser) un quasi-isomorphisme \(E\longrightarrow F\) avec \(E\) strictement \(n\)-pseudo-cohérent et \(E^i=0\) pour \(i\geq 1\), i.e. une \(n\)-présentation finie de \(F\). La deuxième assertion se démontre de manière analogue (utiliser (2.7.1)).

\textbf{Proposition 2.10.} \textbf{a)} \emph{Soient \(m,n\in\mathbb Z\), avec \(n\leq m\). Si \(F\in\Ob\Dpx^-(\mathcal C_S)\) est \(n\)-pseudo-cohérent et acyclique en degré \(\geq m\), il existe localement un quasi-isomorphisme \(E\longrightarrow F\), avec \(E\) strictement \(n\)-pseudo-cohérent et \(E^i=0\) pour \(i\geq m\). Si \(S\) est ponctuel, et \(F\) pseudo-cohérent et acyclique en degré \(\geq m\), il existe un quasi-isomorphisme \(E\longrightarrow F\) avec \(E\) strictement pseudo-cohérent et \(E^i=0\) pour \(i\geq m\).}

\textbf{b)} \emph{Soit \(F\in\Ob\Dpx^-(\mathcal C_S)\) acyclique en degré \(>n\). Pour que \(F\) soit \(n\)-pseudo-cohérent, il faut et il suffit que \(\Ho^n(F)\) soit de type fini.}

\emph{Preuve.} a) Comme \(F\) est acyclique en degré \(\geq m\), la flèche \(0\to F\) est un \(m\)-quasi-isomorphisme. L'assertion résulte de (2.4) et (2.7.1).

b) Il existe un triangle distingué de \(\Dpx^-(\mathcal C_S)\) :
\[
\begin{tikzcd}
& F' \arrow[dl,"+1"'] & \\
F \arrow[rr] && \Ho^n(F)[-n] \arrow[ul] & ,
\end{tikzcd}
\]
(pour le voir, commencer par tronquer \(F\) en remplaçant \(F^n\) par \(Z^n(F)\) et \(F^i\) par \(0\) pour \(i>n\)). La suite exacte de cohomologie implique \(\Ho^i(F')=0\) pour \(i\geq n-1\), donc \(F'\) est \((n-1)\)-pseudo-cohérent. On conclut grâce à (2.5) et (2.6).



% Additional notation for Expose I continuation
\providecommand{\RHom}{\mathrm R\!\Hom}
\providecommand{\Ext}{\operatorname{Ext}}
\providecommand{\coh}{\operatorname{coh}}
\providecommand{\PreCoh}{\operatorname{precoh}}
\providecommand{\Tr}{\operatorname{Tr}}
\providecommand{\Qis}{\operatorname{Qis}}
\providecommand{\Rad}{\operatorname{rad}}
\providecommand{\Spec}{\operatorname{Spec}}
\providecommand{\End}{\operatorname{End}}
\providecommand{\Id}{\operatorname{Id}}
\providecommand{\varinjlim}{\mathop{\mathrm{lim}}\limits_{\longrightarrow}}
\providecommand{\varprojlim}{\mathop{\mathrm{lim}}\limits_{\longleftarrow}}
\providecommand{\qcoh}{\operatorname{qcoh}}
\providecommand{\amp}{\operatorname{amp}}


